\clearpage
\section{Inter-rater Agreement} \label{app:validation}

To assess the reliability of our annotation procedure and the consistency of coding judgments across reviewers, we measured inter-rater agreement on a randomly selected subset of 46 benchmark papers. Each paper in this subset was independently annotated by two reviewers using our codebook of 30 categorical items relating to phenomena, tasks, metrics, and validity claims.

Inter-rater agreement is essential in systematic reviews where subjectivity or interpretation may influence labelling decisions. High agreement indicates that the annotation schema is sufficiently well-defined and that results can be considered reliable across different coders. Conversely, low agreement may indicate ambiguity in the coding process. Prior work in empirical machine learning has emphasized the importance of such reliability assessments when coding qualitative features or design properties \cite{artstein2008inter, geiger2020garbage}.

Given the mix of binary, multi-class, and multi-label questions in our codebook, and the presence of strong label imbalance in many fields, we report agreement using both \textit{Percent Agreement} and \textit{Brennan–Prediger Kappa (BPK)} \cite{bpkappa}.

Percent agreement is computed as the proportion of items on which both raters agreed. For binary and multi-class fields, agreement is determined via exact label match. For multi-label questions (i.e. `check all that apply'), we compute the Jaccard similarity between the two sets of selected options for each item and take the average of these values across all items as the percent agreement score.

BPK adjusts for chance agreement under the assumption of uniform response distributions and is more robust to label imbalance than standard chance-corrected metrics such as Cohen’s Kappa, Fleiss’ Kappa, or Krippendorff’s Alpha, which often degrade in skewed settings. It is defined as:
\[
\text{BPK} = \dfrac{P_o - k^{-1}}{1 - k^{-1}},
\]
where \( P_o \) is the observed proportion of agreement and \( k \) is the number of valid response categories for the question. For binary fields (\(k = 2\)), this simplifies to $\text{BPK} = 2P_o - 1$. 

We compute BPK for all binary and multi-class single-label questions using the corresponding value of \(k\). For multi-label fields, we threshold the Jaccard similarity at 0.3 and treat item pairs with similarity above this threshold as “agreed”, thus converting the comparison into a binary decision before applying BPK. While BPK does not handle missing values, Krippendorff’s Alpha does, it is less suited to class-imbalanced data, and therefore we prefer BPK in our setting.

Across the 30 annotated fields, we observe a mean percent agreement of $68.1\%$ and a mean BPK of $0.524$, indicating overall moderate consistency. Structured and objective fields  exhibit very high agreement. In contrast, interpretive or compositional fields such as \texttt{task\_ecology} and \texttt{dataset\_sampling\_method} show lower consistency, highlighting areas where definitions could be improved.

\begin{table}[h!]
\centering
\caption{Inter-rater agreement across all fields.}
\label{tab:irr-results}
\begin{tabular}{lccc}
\toprule
\textbf{Codebook Field} & \textbf{Question Type} & \textbf{Percent Agreement (\%)} & \textbf{BPK} \\
\midrule
task\_face\_validity          & Multi-class     & 95.12 & 0.927 \\
metric\_access                & Binary          & 95.12 & 0.902 \\
benchmark\_availability       & Multi-class     & 95.12 & 0.939 \\
inclusion                     & Binary          & 93.48 & 0.870 \\
result\_interpretation        & Binary          & 92.68 & 0.854 \\
metric\_aggregation           & Multi-label     & 87.20 & 0.756 \\
metric\_face\_validity        & Multi-class     & 85.37 & 0.780 \\
task\_train\_val              & Multi-label     & 84.15 & 0.951 \\
metric\_metascoring           & Multi-class     & 82.93 & 0.787 \\
results\_human\_baseline      & Binary          & 80.49 & 0.610 \\
task\_dataset\_metadata       & Binary          & 75.61 & 0.512 \\
metric\_fewshot               & Binary          & 73.17 & 0.463 \\
authorship                    & Multi-class     & 73.17 & 0.642 \\
response\_format              & Multi-label     & 71.54 & 0.707 \\
metric\_subscores             & Binary          & 70.73 & 0.415 \\
definition\_integrity\_detail & Multi-class     & 60.98 & 0.415 \\
results\_comparison           & Binary          & 60.98 & 0.220 \\
phenomenon\_short             & Multi-label     & 60.98 & 0.220 \\
results\_realism              & Multi-class     & 58.54 & 0.447 \\
definition\_integrity         & Multi-class     & 58.54 & 0.378 \\
definition\_scope             & Multi-class     & 56.10 & 0.341 \\
results\_comparison\_explanation & Multi-class & 56.10 & 0.341 \\
metric\_definition            & Multi-label     & 55.37 & 0.512 \\
task\_source                  & Multi-label     & 54.23 & 0.561 \\
results\_author\_validity     & Multi-class     & 53.66 & 0.305 \\
phenomenon\_defined           & Binary          & 53.66 & 0.073 \\
phenomenon\_contested         & Multi-class     & 48.78 & 0.317 \\
exclusion\_criteria           & Multi-class     & 40.00 & 0.200 \\
dataset\_sampling\_method     & Multi-label     & 37.80 & 0.122 \\
task\_ecology                 & Multi-class     & 31.71 & 0.146 \\
\midrule
\textbf{Mean}                 & —               & \textbf{68.11} & \textbf{0.524} \\
\bottomrule
\end{tabular}
\end{table}

These findings support the overall reliability of our coding process and offer a basis for targeted improvements in future annotation protocols, especially in interpretive or multi-choice fields.


